\documentclass[11pt]{article}
\usepackage[margin=0.9in]{geometry}
\usepackage{amsmath,amssymb}
\usepackage{enumitem}
\usepackage{tcolorbox}
\usepackage{xcolor}
\usepackage{booktabs}
\usepackage{array}
\usepackage{fancyhdr}
\tcbuselibrary{skins,breakable}

\definecolor{boxbg}{gray}{0.94}
\definecolor{boxline}{gray}{0.55}

\pagestyle{fancy}
\fancyhf{}
\renewcommand{\headrulewidth}{0pt}
\fancyfoot[C]{\small\color{gray} Linear Algebra \quad$\cdot$\quad Pivoting the Simplex Tableau \quad$\cdot$\quad Flashcards}
\fancyfoot[R]{\small\color{gray}\thepage}

\newtcolorbox{cardbox}{enhanced, breakable, colback=boxbg, colframe=boxline,
  boxrule=0.6pt, arc=3pt, left=10pt, right=10pt, top=8pt, bottom=8pt}

\begin{document}
\begin{center}
{\LARGE\bfseries Pivoting the Simplex Tableau}\\[3pt]
\rule{2.2in}{0.8pt}\\[5pt]
{\itshape row operations, pivots, and the simplex decision rules \quad$\bullet$\quad Linear Algebra}
\end{center}

\section*{Quiz Yourself}

Cover the answers page. Say each answer out loud or write it down, then check yourself.

\begin{cardbox}
\textbf{Card 1.} What are the three elementary row operations you are allowed to perform on a matrix?
\end{cardbox}

\begin{cardbox}
\textbf{Card 2.} In a matrix that has been fully reduced, which variables are called \emph{basic} and which are called \emph{free}?
\end{cardbox}

\begin{cardbox}
\textbf{Card 3.} In a simplex tableau for a maximization problem, what rule tells you which variable enters the basis (that is, which column becomes the pivot column)?
\end{cardbox}

\begin{cardbox}
\textbf{Card 4.} Once the pivot column is chosen in a simplex tableau, how do you decide which row to pivot in?
\end{cardbox}

\begin{cardbox}
\textbf{Card 5.} After a pivot, the bottom (objective) row of your maximization tableau has no negative entries in the decision-variable columns. What does this result tell you, and what do you do next?
\end{cardbox}

\newpage

\section*{Answers}

\begin{cardbox}
\textbf{Card 1.} (1) Switch two rows. (2) Multiply a row by a nonzero number. (3) Replace a row by itself plus a multiple of another row.
\end{cardbox}

\begin{cardbox}
\textbf{Card 2.} Variables whose columns contain a pivot position are \emph{basic} variables; variables whose columns contain no pivot position are \emph{free} variables. The basic variables can be solved for in terms of the free ones.
\end{cardbox}

\begin{cardbox}
\textbf{Card 3.} Look at the objective (bottom) row and choose the column with the \emph{most negative} entry. That column's variable enters the basis, because increasing it improves the objective value fastest.
\end{cardbox}

\begin{cardbox}
\textbf{Card 4.} Use the minimum-ratio test: for every row whose entry in the pivot column is positive, divide the right-hand-side value by that entry, and pivot in the row with the \emph{smallest} ratio. For example, if the pivot column has entries $5$ and $2$ with right-hand sides $40$ and $30$, the ratios are $40/5 = 8$ and $30/2 = 15$, so you pivot in the first row.
\end{cardbox}

\begin{cardbox}
\textbf{Card 5.} The tableau is \emph{optimal}: no variable can enter the basis and increase the objective any further. You stop pivoting and read the maximum objective value from the right-hand side of the bottom row.
\end{cardbox}

\vfill

{\small\color{gray}
``A First Course in Linear Algebra'' by Ken Kuttler, used under CC BY 4.0.\\
``Understanding Linear Algebra'' by David Austin, used under CC BY 4.0.
}

% SOURCES
% Card 1: #kuttler-1.3-def-2 (Definition 1.3.2, Elementary Row Operations)
% Card 2: #principle-1-4-1 (Austin, Principle 1.4.1, basic vs. free variables)
% Card 3: lo_mapping:missing_concepts[0] (entering-variable rule)
% Card 4: lo_mapping:missing_concepts[1] (minimum-ratio test)
% Card 5: lo_mapping:missing_concepts[2] (simplex optimality criterion)
\end{document}
